\section{Introduction}
\label{sec:intro}

Large language models (LLMs) are now trained and deployed across a rapidly
diversifying hardware landscape.  While NVIDIA GPUs and the CUDA software stack
have historically dominated LLM development, AMD Instinct accelerators running
ROCm are increasingly available in data centers and cloud offerings.  Yet most
published training results are implicitly CUDA-centric, leaving practitioners
uncertain about whether their algorithms, training frameworks, and inference
engines work correctly and efficiently on non-NVIDIA hardware.

Answering this question rigorously requires a \emph{controlled} experimental
vehicle: a codebase that implements every algorithm from scratch in portable
PyTorch, with no hidden vendor-specific dependencies, so that any observed
performance difference is attributable to the hardware and runtime stack rather
than to library porting gaps.

\paragraph{The MiniMind vehicle.}
MiniMind~\citep{minimind2024} is an end-to-end LLM pipeline implemented in pure
PyTorch without relying on \texttt{trl}, \texttt{peft}, or
\texttt{transformers.Trainer}.  It covers the full training lifecycle in a
single, coherent codebase: BPE tokenization, causal LM pretraining, supervised
fine-tuning (SFT), LoRA adaptation, knowledge distillation, direct preference
optimization (DPO), proximal policy optimization (PPO), group-relative policy
optimization (GRPO), and agent reinforcement learning with tool-call rollouts.
Its \texttt{DeviceCtx} abstraction~(\S\ref{sec:methodology}) routes all
device-specific calls — mixed-precision autocast, gradient scaling, distributed
rank binding — through a single layer that works unchanged on both CUDA and ROCm,
making it possible to run the \emph{identical} trainer scripts on both platforms.

\paragraph{Contributions.}
This paper makes the following contributions:
\begin{enumerate}[leftmargin=*]
  \item \textbf{Full-pipeline H100 vs MI350 comparison.}  We run every
    MiniMind training stage on single-node 1-, 4-, and 8-GPU configurations of
    both NVIDIA H100 (Hopper, CUDA~12, NCCL) and AMD Instinct MI350 (CDNA4,
    ROCm~7.1, RCCL), reporting throughput, model-FLOP utilization (MFU), peak
    HBM consumption, and convergence parity for each.

  \item \textbf{Dense and MoE coverage.}  We benchmark both the dense
    ($\approx$64\,M-parameter) and mixture-of-experts
    ($\approx$198\,M-A64\,M) model variants, explicitly exercising
    MoE-specific token-dispatch and router load-balancing kernels on each vendor.

  \item \textbf{Inference-stack benchmarks.}  Using the in-repo Qwen3
    conversion path, we benchmark HF-native \texttt{generate()}, vLLM, and
    SGLang on both platforms across standard workloads: latency, throughput,
    chain-of-thought reasoning (\texttt{<think>} decoding), and tool-call
    evaluation.

  \item \textbf{Cross-cutting analysis.}  We characterize convergence parity,
    \texttt{torch.compile} codegen quality, dtype maturity (bf16 / fp16 / FP8
    outlook), MoE routing behavior, and intra-node collective communication
    bandwidth (NCCL vs.\ RCCL) as cross-cutting dimensions.

  \item \textbf{A reproducible benchmark artifact.}  We release all
    experiment scripts, metric-collection utilities, and exact software version
    pins so that the community can re-run and extend these results as hardware
    and software stacks evolve.
\end{enumerate}

\paragraph{Paper organization.}
Section~\ref{sec:background} describes the MiniMind architecture, the training
algorithms covered, and the hardware/software stacks under test.
Section~\ref{sec:methodology} details our experimental methodology, including
the vendor-portable code path, reproducibility recipe, and MFU computation.
Sections~\ref{sec:training}--\ref{sec:inference} present training and inference
results respectively.  Section~\ref{sec:crosscutting} provides cross-cutting
analyses.  Sections~\ref{sec:discussion}--\ref{sec:related} discuss implications,
limitations, and related work, and Section~\ref{sec:appendix} contains the full
reproducibility appendix.
